% Source fingerprint: 16e6116a710ba5c91b5a96ada28fde11a82b2798400c45c10503adb761954a6d
% T15: raw TeX cells preserved; no numeric highlighting.
\begin{table}[htbp]
\centering\small\renewcommand{\arraystretch}{1.18}\setlength{\tabcolsep}{4pt}
\caption[有界变差函数的三种收敛]{有界变差函数的三种收敛。在已有全域 \(L^1\) 收敛与统一变差界时，正文可将分布收敛升级为导数测度弱星收敛；不同收敛方式仍有不同的附加条件。}
\label{b1:tab:visual-t15}
\begin{tabularx}{\linewidth}{@{}>{\raggedright\arraybackslash}p{3.1cm}>{\raggedright\arraybackslash}p{4.7cm}>{\raggedright\arraybackslash}X@{}}
\toprule
收敛方式 & 需要控制的数据 & 典型差别 \\
\midrule
有界变差范数收敛 & \(\|u_j-u\|_1+|D(u_j-u)|(\Omega)\to0\) & 推出严格收敛，也推出本书定义的弱星收敛 \\
严格收敛 & \(u_j\to u\) 于 \(L^1\)，且 \(|Du_j|(\Omega)\to|Du|(\Omega)\) & 总变差数值相近，不等于导数测度之差的总变差趋零 \\
\WeakStar{}收敛 & 全域 \(L^1\) 收敛，加上导数测度对 \(C_0\) 的弱星收敛 & 质量可以向边界或无穷远逃逸，未必严格收敛 \\
局部版本 & 只在内部相对紧区域控制 \(L^1\) 与变差 & 不能把局部紧性直接提升为粗糙域上的全局紧性 \\
\bottomrule
\end{tabularx}
\end{table}
